% Previous conclusion preserved (2026-09-26).
% \section{Conclusion}
% We present Tandem, a MoT-based framework that synchronously generates videos and motions. We design a coupled noise schedule and a motion representation to improve the motion quality and reduce the jitter. Our ablation studies validate the effectiveness of our strategies. The comparison results with other motion generation methods demonstrate the convenience, scene-awareness, speed and generalization of our framework. Our method benefits the game development by automatically generating character motions, reducing the labor-intensive Keyframe animations. However, our framework can't generate motions in real time and synthesize fine-grained hand interactions, which limits the further application in the game. We will curate hand-interaction datasets and use distillation to achieve faster and detailed interactions in the future.
% \label{sec:conclusion}
% 
% 


% Previous conclusion preserved before separate Limitations section.
% \section{Conclusion}
% \label{sec:conclusion}
% We present World2Motion, which adapts Cosmos 3 to jointly generate scene-aware 3D human motion and video from an image and a text prompt. A hybrid dataset supports this adaptation, while relative root translation and a shift-decoupled noise schedule improve temporal stability. Experiments show improved motion--text alignment and competitive scene interaction, with approximately $3.34\times$ faster inference than the evaluated two-stage pipeline. Contact accuracy remains a limitation, as the two-stage baseline retains a higher Contact-F1. Generation remains slower than real time, and fine-grained hand interactions remain challenging. Future work includes improving inference efficiency and providing more precise supervision for hand--object contact.

\section{Limitations}
\label{sec:limitations}
% Before clarifying data coverage and evaluation limits: We view World2Motion as an initial baseline for adapting pretrained video world models to 3D human motion generation. Our training data cover a limited range of interactions, and performance on more complex cases has not been systematically evaluated. CameraHMR errors can affect both initial poses and supervision from real videos. Inference is faster than CoMoVi in our comparison, but remains too slow for real-time applications. Reliable root placement and precise contact remain challenges, with Contact-F1 below the two-stage baseline. Richer interaction data, additional modalities such as scene geometry and object states, and faster generation are directions for future work. For current offline applications, generated motions can initialize agent-assisted workflows for inspection and refinement. The supplementary video illustrates such a workflow, including motion smoothing and contact refinement. These examples illustrate a use of our outputs; broader gains in workflow quality and efficiency require systematic evaluation.
% Previous Limitations opening: We view World2Motion as an initial baseline for adapting pretrained video world models to 3D human motion generation.
% Before author-approved rewrite: World2Motion provides a baseline for adapting pretrained video world models to 3D human motion generation. Training data mainly cover single-person motion and simple interactions; more complex cases lack systematic evaluation. CameraHMR errors can affect initial poses and real-video supervision. Inference is faster than CoMoVi in our comparison, but remains unsuitable for real-time applications. Our metrics provide a limited assessment: ISR and PoseError do not guarantee precise contact. Root placement and contact remain challenges, with Contact-F1 below the two-stage baseline. Future directions include richer interactions, additional modalities such as scene geometry and object states, and faster generation. For offline applications, our outputs can initialize agent-assisted inspection and refinement. The supplementary video illustrates motion smoothing and contact refinement in such a workflow; broader quality and efficiency gains require systematic evaluation.
World2Motion provides a baseline for adapting pretrained video world models to 3D human motion generation, but several limitations remain. First, our training data mainly cover single-person motion and simple interactions, leaving more complex scenarios largely unexplored. Second, despite the improvement in inference speed, World2Motion is still too slow for real-time use. Third, accurate root placement and contact remain challenging. Our evaluation also has limits: a high ISR or low PoseError does not imply precise contact. Future work could expand interaction coverage, incorporate scene geometry and object states, and reduce inference time. For now, generated motions can serve as starting points for further editing. The supplementary video shows examples of agent-assisted smoothing and contact refinement.

\section{Conclusion}
\label{sec:conclusion}
% Before removing the repeated baseline-positioning sentence: World2Motion adapts Cosmos 3 to jointly generate scene-aware 3D human motion and video from an image and a text prompt. Our adaptation combines hybrid video--motion training data, relative root translation, and a shift-decoupled noise schedule. The ablations show reduced jerk and foot sliding under the tested settings. This work provides a baseline for further study of video-model adaptation to image-conditioned motion generation.
% Before author-approved rewrite: World2Motion adapts Cosmos 3 to jointly generate scene-aware 3D human motion and video from an image and a text prompt. Our adaptation combines hybrid video--motion training data, relative root translation, and a shift-decoupled noise schedule. The ablations show reduced jerk and foot sliding under the tested settings.
We present World2Motion, which adapts Cosmos 3 to jointly generate scene-aware 3D human motion and video from a single image and a text prompt. We train the model on synthetic video--motion pairs and real videos with estimated motion, and use a shift-decoupled noise schedule to reduce motion jitter. On our interaction benchmark, World2Motion improves motion--text alignment over the compared methods and matches the interaction success rate of the two-stage baseline.
